% 第3章 3.2.2：紧自伴算子沿特征方向的伸缩。
% (a) l^2 上 A=diag(1/j)，u 的系数 <e_j,u>=1/j，Au 的系数 lambda_j<e_j,u>=1/j^2，j=1..10。
% (b) L^2([-1,1],dz/2) 上核 1-6xz：e_1=1 的特征值 1，e_2=sqrt3 x 的特征值 -2
%     （T x = -6x * int z^2 dmu = -2x）。在 (e_1,e_2) 系数平面上单位圆映为半轴 1、2 的椭圆；
%     点 p=(cos45°,sin45°)=(0.7071,0.7071) 映为 Tp=(0.7071,-1.4142)。
\begin{tikzpicture}[font=\figurefont\footnotesize]
\begin{axis}[name=A,width=58mm,height=40mm,
  ybar,bar width=1.4mm,xmin=0.3,xmax=10.7,ymin=0,ymax=1.1,
  axis x line=bottom,axis y line=left,
  xtick={1,2,...,10},ytick={0,0.5,1},
  xlabel={$j$},ylabel={系数},
  xlabel style={font=\figurefont\scriptsize},ylabel style={font=\figurefont\scriptsize},
  tick label style={font=\figurefont\scriptsize,text=FigureInk},
  ymajorgrids,grid style={draw=FigureGrid,line width=.4pt},
  legend style={at={(0.98,0.97)},anchor=north east,font=\figurefont\scriptsize,cells={anchor=west}},
  title={(a) 系数逐项乘以$\lambda_j=1/j$},title style={font=\figurefont\footnotesize}]
\addplot[fill=FigureBlueFill,draw=FigureBlue] coordinates {(1,1) (2,0.5) (3,0.3333) (4,0.25) (5,0.2) (6,0.1667) (7,0.1429) (8,0.125) (9,0.1111) (10,0.1)};
\addplot[fill=FigureBlue,draw=FigureBlue] coordinates {(1,1) (2,0.25) (3,0.1111) (4,0.0625) (5,0.04) (6,0.0278) (7,0.0204) (8,0.0156) (9,0.0123) (10,0.01)};
\legend{{$\langle e_j,u\rangle$},{$\langle e_j,Au\rangle$}}
\end{axis}
\begin{scope}[shift={($(A.east)+(22mm,-1mm)$)},scale=0.78]
  \draw[FigureGrid,line width=.4pt,-{Stealth[length=1.6mm]}] (-1.5,0) -- (1.6,0) node[right,text=FigureInk] {$e_1$};
  \draw[FigureGrid,line width=.4pt,-{Stealth[length=1.6mm]}] (0,-2.25) -- (0,2.3) node[above,text=FigureInk] {$e_2$};
  \fill[FigureBlueFill] (0,0) ellipse (1 and 2);
  \draw[FigureBlue,line width=1.1pt] (0,0) ellipse (1 and 2);
  \draw[FigureStroke,line width=.8pt,dashed] (0,0) circle (1);
  \fill[FigureStroke] (0.7071,0.7071) circle (1.3pt) node[above right=-1pt,text=FigureInk] {$p$};
  \fill[FigureBlue] (0.7071,-1.4142) circle (1.3pt) node[right,text=FigureInk] {$Tp$};
  \draw[FigureMuted,line width=.5pt,-{Stealth[length=1.4mm]}] (0.7071,0.62) -- (0.7071,-1.32);
  \node[text=FigureInk,anchor=south west,font=\scriptsize] at (1.0,0.1) {$\lambda_1=1$};
  \node[text=FigureInk,anchor=west,font=\scriptsize] at (0.1,1.9) {$|\lambda_2|=2$};
  \node[anchor=north] at (0.05,-2.45) {(b) 负特征值使方向翻转};
\end{scope}
\end{tikzpicture}
